\section{Which counterfactual questions have enough evidence?}

A rare exposure can carry a large inverse-probability weight, and shared
randomness can make many observed rows convey less information than their count
suggests. These two difficulties interact with the question being asked.
A policy change supported on a small region should not automatically pay for
an unrelated rare exposure elsewhere; a subgroup contrast should not require
two imprecise policy levels when its signed residual cancels outside that
subgroup. The relevant object is therefore a family of explicit, supported
counterfactual moments together with their evidence requirements.

We formulate these questions conditional on an exogenous design. Rows can
have different conditional means and can form one connected dependency graph.
The graph is sufficient for the entire random score vector, with independence
between disjoint sets that have no crossing edge. This requirement is stronger
than a drawing of the application network, pairwise zero covariance, or a graph
inferred from one realized log. Independent random-source incidence provides
a constructive sufficient model.

The main construction separates the known direct prediction term from the
random weighted residual. It clips the residual at a threshold fixed before
certification outcomes and pays for both the resulting worst-case bias and
graph-adjusted fluctuations. The resulting interval covers the original
unclipped signed moment. The statistical ingredients are standard clipping
and concentration tools; graph-coloring concentration has a substantial prior
literature~\cite{janson2004}. We formulate and analyze how support, contrasts,
residual moments and a sufficient graph determine the evidence cost of an
auditing question. A generic truncated graph mean with fully matched
information is the same procedure; the contribution therefore concerns this
formulation and analysis, its limits and its decision consequences.

Independent construction followed by risk-controlled selection has clear
precedents. Learn then Test combines valid tests with familywise error control
to protect a returned set~\cite{angelopoulos2025ltt}; high-confidence policy
improvement separates candidate search from a held-out safety
test~\cite{thomas2015hcpi}. In the present model, this logic requires valid
certificates for heterogeneous signed moments under a sufficient conditional
graph. The analysis connects those certificates to support, residual moments,
paired comparisons and a scoped evidence-cost obstruction.

Three questions must be kept separate. First, are the assumptions sufficient
and the implementation consistent with the stated theorem? Second, is the
resulting interval informative for the given subgroup, overlap and graph?
Third, does it return a policy that is simultaneously safe and useful?
The historical confirmation study illustrates why coverage alone cannot answer
the latter two questions. The experiments evaluate these distinctions under
heterogeneous connected graphs and independent discovery. The completed
validity grid already shows that full-family coverage can coexist with zero
detection in every tested heavy-tail cell, while stronger bounded inference
can recover useful detection. Query-local ablations and a fully informed
equal comparator separate gains from relevant information from gains
attributable to a distinct statistical procedure.
